\section{Bounding the denominator sum with coarser local estimates}

% BEGIN E83D1
\paragraph{E83D1. A signed excess on the original congruence quotient.}
Work on a late qualifying family member and its common hub, with
the hypotheses of G82I1--G82I2. For a guarded inert C or D index
$i$ and a good point $z$, use $p=\ell_i$, $n_V=r+t$,
$v=\operatorname{vdel}(i)$, and $m_p=\#\{u:u\mid p\}$ in
the CM host $E$. Put $d_u=[E_u:\mathbb Q_p]$ and
$b_u=X.\mathrm{lenRow}(i,u)$. Write $A_{iz}$ for the complete
right hand side of the required pointwise bound, including the
terms outside the row:
\[
 A_{iz}=s_z+\sum_{0<a<t}u_a+v_p(h_E)+p^{-2}
     +\sum_{u\mid p}b_u+tc_B\sum_{u\mid p}(d_u+2)+4n_Vv .
\]
Here $u_a$ denotes a cyclotomic length. For each $u$ take any
proved nonnegative upper bound $B_{izu}$ for the normalized
cokernel of every tested dual row map there. Set
\[
 D_{iz}=\sum_{u\mid p}(B_{izu}-b_u)-tc_Bm_p-(2n_V+2)v.
\]
Then $v_{\rm cc}(i,z)\le A_{iz}+D_{iz}$. At a point where the
native row inequality has already been proved, $D_{iz}$ may be
replaced by $\min(0,D_{iz})$. This follows by substituting the
actual row and character terms from G82I1 in G82I2 and
subtracting the two expressions, before taking any maximum.
At every remaining point N82I2 supplies, without a height guard
on this estimate itself,
\[
 B_{izu}=(d_u+1)r(25h_B/p+3).
\]
Smaller proved bounds can be inserted place by place. Since the
host has degree $2f$, $\sum d_u=2f$ and $m_p\le2f$.
Nonnegativity of the native allowances therefore gives the
uniform alternative $D_{iz}^+\le4fr(25h_B/p+3)$, where
$a^+=\max(0,a)$ and the displayed coarse choice of $B$ is used.
% END E83D1

% BEGIN E83D2
\paragraph{E83D2. Multiplicity and place weights for a pointwise excess.}
Let $j$ be an allowed lift, $d_0=d_{j,0}$, and write $w$ for a
place of the coefficient field $K$ (distinct from the host
places $u$). In the native sum only twisted places in the finite
support occur. Write $e_w,f_w$ for the ramification and residue
degrees above the rational prime $p$, and
$\omega_w=e_w f_w\log p/[K:\mathbb Q]$. At an ordinary prime
outside the two generator primes use the following finite sum:
\[
 W_{jp}(a)=\sum_{w\in\mathrm{twAt}(p)}\omega_w
 \sum_{i\in I_{jw}}\sum_{z\in Z_{\rm good}(i)}m_{iz}a_{iz},
 \qquad m_{iz}=\dim V_{iz}.
\]
If $v_{\rm cc}(i,z)\le A_{iz}+D_{iz}$, the contribution at this
prime is at most
\[
       W_{jp}(A+v_p(M_0))+W_{jp}(D),\qquad
       M_0=\mathrm{restRowsFiniteGroupCard}.
\]
For point types already known to satisfy their bound use their
own prescribed $A$ and $D=0$. To see the inequality, the
$\mathrm{pass}$ field bounds the normalized length of the
good projection quotient at $i$ by
$\sum_z m_{iz}(v_{\rm cc}(i,z)+v_p(M_0))$. This is the
point at which the difference between the Hecke support $Z$
and its good subset matters. LN3--LN6 give the lattice and the
auxiliary projector cost. The tie for $\eta_j(w)$ sums exactly
these lengths over $I_{jw}$. Multiply by $\omega_w$ and add.
Thus a signed improvement obtained at one good point survives
with its actual multiplicity. No division by $d_0$ occurs.

For a bound uniform over the points at $p$, replace the
second term by $d_0\log p\,\max(0,\max D_{iz})$, taking zero
for an empty set. Indeed $\sum_{i,z}m_{iz}\le d_0$ at each
such $w$, by $\mathrm{mult}$ and $\mathrm{rank}$ on the same
member. The places in the support form a subset of the primes
of $K$ over $p$, so $\sum_w e_wf_w\le[K:\mathbb Q]$. This
also proves the assertion at $d_0=0$. The separate generator
prime estimate is used in its own form and acquires no excess
from these ordinary middle-range rows.
% END E83D2

% BEGIN E83D3
\paragraph{E83D3. The entire exceptional-prime excess fits the global budget.}
Let $p_*=\ell_*(L)$, and use the exceptional set $E_L$ of X78A2.
For every allowed $j$, the excess in E83D2 coming from all
C/D points estimated here only by the coarse bound at $p>p_*$ is bounded by
\[
       d_0\,27648\,U^6,\qquad
       U=x^{C+2}\exp(x^{\epsilon/4}).
\]
Here $C\ge1$ is a fixed common family parameter exponent.
In particular this bound is $d_0\,o(b)$ uniformly over the
members, hubs and lifts, where $b=\exp(x^\epsilon)$.

Indeed a middle-range prime above $p_*$ belongs to $E_L$ by the
definition of its guard. X78A1 gives $12b/p_*\le\eta/m$.
After one family cutoff $h_B\le2b$. Thus at $p>p_*$,
$h_B/p\le\eta/(6m)$. We record the numerical bounds for $U$.
A78P3 gives $n,k,t\le x^C$ on increasing $C$ once. The
parameter inequalities $\log\eta\le x^{\epsilon/4}$,
$\log\Delta\le x^{\epsilon/4}$ and $S_c\le4k\Delta$
give $\eta\le U$ and $S_c\le4U$.
Also $f=2n\le2U$, $r\le2k\le2U$,
$1\le x^\epsilon\le U$, and
$|\mathrm{TT}|\le N_\star\le2n\le2U$.
With these bounds $h_B=b+S_c+\eta+t+2\le2b$ holds
uniformly since $U/b\to0$.

We now count the relevant primes without using a distribution
theorem. The upper bound in X78A2 is
\[
 (|\mathrm{TT}|+2)(S_c+3\eta+t+8b/p_*+9)\le72U^2 .
\]
The first factor is at most $4U$. In the second,
$8b/p_*\le(2/3)\eta/m\le U$ and the five summands are at
most $4U,3U,U,U,9U$. Every exceptional prime is at most $3b$,
so $\log p\le3U$. E83D1 bounds the positive excess by
$4fr(25U/6+3)\le128U^3$. E83D2 now bounds the loss over the
entire set by $d_0(72U^2)(3U)(128U^3)$, as claimed.
Finally
\[
  \log(U^6/b)=6(C+2)\log x+6x^{\epsilon/4}-x^\epsilon
                \longrightarrow-\infty.
\]
The cutoff for any fixed desired fraction of $b$ precedes every
member and point. In the global family estimate this entire
additional term costs at most one extra unit of the constant
eventually, since $fm(x^\epsilon)^2\ge1$. This controls the
error relative to the prescribed right hand sides at the
exceptional primes. Their Selmer, cyclotomic and baseline
contributions still enter the usual denominator sum. This aggregate estimate can be used independently of the pointwise
argument. A83G1--A83G2 supply the pointwise inequality even at
the exceptional primes in its original range.
% END E83D3
